% ============================================================================
% 2.6 主程序：把一切串起来
% ============================================================================
\subsection{hello.s：主程序入口}

主程序的工作非常简单——初始化堆之后，就进入一个循环：
读一行→复制→存到平行数组→打印→找空格。

\begin{lstlisting}[style=asm,caption=hello.s — 主程序]
.text
.global _main
.extern exit
.extern memory_init
.extern memory_alloc
.extern memory_free
.extern sys_write
.extern sys_read
.extern test_write
.extern test_memory
.extern test_read
.extern test_write_memory
.extern store_string_ptr_len
.extern get_string_ptr
.extern get_string_len
.extern store_ptr_done
.extern print_index_strings
.extern find_two_spaces

_main:
    bl memory_init             ; 初始化堆
    mov x19, #0               ; x19 = 循环计数器（也是字符串序号）

loop:
    mov x20, #0
    bl store_string           ; 读一行，分配内存，存入数组

    mov x20, #0
    bl get_string_ptr_len_spaces ; 顺便检查空格位置

    add x19, x19, #1
    cmp x19, #1
    b.eq do_exit
    b loop

do_exit:
    b exit
\end{lstlisting}

\subsection{store\_string：读入并保存一行}

\begin{lstlisting}[style=asm,caption=hello.s — store\_string]
store_string:
    stp x29, x30, [sp, #-16]!
    bl test_read              ; 把一行
    bl copy_string            ; copy_string(x0) 从 input_buf 复制到新分配的堆内存
    mov x1, x20                ; 序号
    bl store_string_ptr_len     ; 存入平行数组
    ldp x29, x30, [sp], #16
    ret
\end{lstlisting}

\subsection{get\_string\_ptr\_len\_spaces：打印并解析空格}

\begin{lstlisting}[style=asm,caption=hello.s — 打印+解析]
get_string_ptr_len_spaces:
    stp x29, x30, [sp, #-16]!

    bl print_string           ; 把刚存入的字符串再取出来打印
    stp x1, x2, [sp, #-16]! ; 保存打印返回的指针和长度
    bl find_two_spaces        ; 然后扫描字符串找空格位置
    ldp x2, x3, [sp], #16

get_string_ptr_len_spaces_done:
    ldp x29, x30, [sp], #16
    ret
\end{lstlisting}

\subsection{理解整体数据流}

整个程序的数据流如下：

\begin{verbatim}
stdin → sys_read → input_buf → copy_string(从 input_buf 到 heap)
        → (ptr, len) → store_string_ptr_len
        → print_index_strings(x19) → sys_write → stdout
        → find_two_spaces(buffer) → x0=第一个空格, x1=第二个空格
\end{verbatim}

程序循环一次（\texttt{x19} 递增一，直到 \texttt{x19 = 1} 时退出。
你可以把 \texttt{cmp x19, \#1} 改成更大的数字来延长程序执行更多轮。

\subsection{主循环限制}

这里的循环在 \texttt{x19} 计数，
\textbf{最多能存 MAX\_STRINGS（32）条字符串}。超过会被
\texttt{store\_string\_ptr\_len} 里的 \texttt{cmp x1, \#MAX\_STRINGS; b.hs store\_ptr\_done} 拦截掉——如果 \texttt{x19} 超过 32，函数什么都不做就返回。

\subsection{新版主程序：把解析后的 token 逐条打印}

在引入 \texttt{store\_instr\_addrs} 之后，主程序多了一步：
把刚才那条命令的 4 个 token 分别打印出来，验证解析是否正确。

\begin{lstlisting}[style=asm, caption=hello.s — 新版 \_main（核心循环）]
_main:
    bl memory_init                 ; 初始化堆
    mov x19, #0                    ; x19 = 循环计数器

loop:
    mov x20, #0
    bl  store_string               ; 读一行 → 复制 → 存到平行数组

    mov x20, #0
    bl  get_string_ptr_len_spaces  ; 顺便找两个空格位置，顺便把字符串打印出来
                                   ; 返回后：x0 = 第一个空格, x1 = 第二个空格
                                   ; x2 = 字符串起始, x3 = 字符串长度

    bl  store_instr_addrs          ; 把刚才的一行切成 3 个 token，存进 instr_addrs

    ; ---- 把 token 0/1/2/3 分别打印出来，每个后面换行 ----
    mov x0, #0
    bl  get_instr_addrs            ; x0 = 第 0 条指令的 4 指针数组头
    bl  get_instr_0                ; 打印 token 0
    bl  print_newline

    mov x0, #0
    bl  get_instr_addrs
    bl  get_instr_1                ; 打印 token 1
    bl  print_newline

    mov x0, #0
    bl  get_instr_addrs
    bl  get_instr_2                ; 打印 token 2
    bl  print_newline

    mov x0, #0
    bl  get_instr_addrs
    bl  get_instr_3                ; 打印 token 3
    bl  print_newline

flag:
    add x19, x19, #1
    cmp x19, #1
    b.eq do_exit
    b loop

do_exit:
    b exit
\end{lstlisting}

调用完空格查找函数后，\texttt{x0/x1} 里是两个空格位置，
\texttt{x2/x3} 里是字符串的起始地址和长度，
正好匹配 \texttt{store\_instr\_addrs} 期望的 4 个入参。

\subsection{get\_instr\_0/1/2/3：从指针数组里拿一个 token 打印}

这四个函数的结构几乎一样，只是读取的偏移不同：

\begin{lstlisting}[style=asm, caption=hello.s — get\_instr\_0（打印 token 0）]
get_instr_0:
    stp x29, x30, [sp, #-32]!
    stp x19, x20, [sp, #16]
    mov x20, x0                    ; 保存 4 指针数组头

    ; 手动 strlen: 找到 token 0 的长度
    ldr x1, [x20]                  ; x1 = token 0 的字符串起始
    mov x2, #0
strlen_loop:
    ldrb w3, [x1, x2]
    cbz  w3, strlen_done
    add  x2, x2, #1
    b    strlen_loop
strlen_done:
    mov  x0, #1
    bl   sys_write                 ; sys_write(1, token0_ptr, token0_len)

    ldp x19, x20, [sp, #16]
    ldp x29, x30, [sp], #32
    ret
\end{lstlisting}

\texttt{get\_instr\_1/2/3} 的区别只是把 \texttt{ldr x1, [x20]}
换成 \texttt{ldr x1, [x20, \#8]}、\texttt{ldr x1, [x20, \#16]}、
\texttt{ldr x1, [x20, \#24]}——分别取 token 1、2、3 的指针。

\subsection{print\_newline：输出一个换行符}

在两个 token 之间打印 \texttt{\textbackslash n}，让每个 token 单独占一行：

\begin{lstlisting}[style=asm, caption=hello.s — print\_newline]
print_newline:
    stp  x29, x30, [sp, #-32]!
    strb wzr, [sp, #16]            ; 清一块栈内存
    mov  w8, #10                   ; '\n' 的 ASCII = 10
    strb w8,  [sp, #16]            ; 把换行符写进去
    add  x1,  sp, #16              ; x1 = 换行符所在地址
    mov  x0,  #1
    mov  x2,  #1
    bl   sys_write                 ; sys_write(1, "\n", 1)
    ldp  x29, x30, [sp], #32
    ret
\end{lstlisting}

它用栈上的一小块空间临时存放换行符，不需要专门准备一个 \texttt{.data}
里的字符串。这是"不依赖静态数据、临时在栈上构造参数"的常见小技巧。

\subsection{现在整个程序的工作流}

综合起来，一次循环做了以下 4 件事：

\begin{enumerate}
    \item 从 \code{stdin} 读一行 → \code{copy\_string} 复制到堆 →
        把指针和长度存入 \code{string\_ptrs} 与 \code{string\_lens} 两个数组
    \item 打印那一行 → 扫描字符串找两个空格位置
    \item 用空格位置把一行切成 3 个 token → 存入 \code{instr\_addrs}
    \item 按序号把 4 个 token 分别打印出来
\end{enumerate}

这相当于一个最小的"命令解析 + 回显"循环：
你输入什么，它就原样打印，再把每个 token 拆出来单独打印一行。
